% [12pt,a4paper,bibliography=totocnumbered,listof=totocnumbered]{scrartcl}
\documentclass[12pt,a4paper]{article}

\usepackage[ngerman]{babel}
\usepackage[utf8]{inputenc}
\usepackage{amsmath}
\usepackage{amsfonts}
\usepackage{amssymb}
\usepackage{graphicx}
\usepackage{fancyhdr}
\usepackage{tabularx}
\usepackage{geometry}
\usepackage{setspace}
\usepackage[right]{eurosym}
\usepackage[printonlyused]{acronym}
\usepackage{subfig}
\usepackage{floatflt}
\usepackage{float}
\usepackage[usenames,dvipsnames]{color}
\usepackage{colortbl}
\usepackage{paralist}
\usepackage{array}
\usepackage{titlesec}
\usepackage{parskip}
\usepackage[right]{eurosym}
%\usepackage{picins}
\usepackage[subfigure,titles]{tocloft}
\usepackage[pdfpagelabels=true]{hyperref}

\usepackage{listings}

%\usepackage[
%backend=biber,
%style=alphabetic-verb,
%citestyle=alphabetic-verb
%]{biblatex}
%\usepackage{biblatex} 
%\addbibresource{literatur.bib}


\lstset{basicstyle=\ttfamily\tiny, captionpos=b, breaklines=true, showstringspaces=false, tabsize=2, xleftmargin=0em, framexleftmargin=0em}
\makeatletter
\def\l@lstlisting#1#2{\@dottedtocline{1}{0em}{1em}{\hspace{1,5em} Lst. #1}{#2}}
\makeatother

\geometry{a4paper, top=27mm, left=20mm, right=20mm, bottom=35mm, headsep=10mm, footskip=12mm}


\hypersetup{unicode=false, pdftoolbar=true, pdfmenubar=true, pdffitwindow=false, pdfstartview={FitH},
	pdftitle={Wahlpflichtfach: Implementierung von Brettspielen am Beispiel ReversiXT (SS \the\year)},
	pdfauthor={Dr.\ Carsten Kern},
	pdfsubject={Projektbericht},
	pdfcreator={\LaTeX\ with package \flqq hyperref\frqq},
	pdfproducer={pdfTeX \the\pdftexversion.\pdftexrevision},
	pdfkeywords={Projektbericht, ReversiXT},
	pdfnewwindow=true,
	colorlinks=true,linkcolor=black,citecolor=black,filecolor=magenta,urlcolor=black}
\pdfinfo{/CreationDate (D:20151500000000)}

\begin{document}

\titlespacing{\section}{0pt}{12pt plus 4pt minus 2pt}{-6pt plus 2pt minus 2pt}

% Kopf- und Fusszeile
\renewcommand{\sectionmark}[1]{\markright{#1}}
\renewcommand{\leftmark}{\rightmark}
\pagestyle{fancy}
\lhead{}
\chead{}
\rhead{\thesection\space\contentsname}
\lfoot{Implementierung von Brettspielen am Beispiel ReversiXT -- SS \the\year}
\cfoot{}
\rfoot{\ \linebreak Seite \thepage}
\renewcommand{\headrulewidth}{0.4pt}
\renewcommand{\footrulewidth}{0.4pt}

% Vorspann
\renewcommand{\thesection}{\Roman{section}}
\renewcommand{\theHsection}{\Roman{section}}
\pagenumbering{Roman}

\newcommand{\folgen}[1]{
\ensuremath
#1
}

\newcommand{\MyTitlepage}[5][\empty]{
\thispagestyle{empty}
\begin{center}
	\includegraphics[scale=0.2]{pics/oth-logo.png}\\
	\vspace*{2cm}
	\Large
	\textbf{Fakultät}\\
	\textbf{Informatik und Mathematik}\\
	\vspace*{2cm}
	\Huge
	\textbf{Projektbericht}\\
	\vspace*{0.5cm}
	\large
	zum HSP1 im Wintersemester 2019/20\\
	\vspace*{1cm}
	\textbf{Implementierung von Reversi mit dem AlphaZero-Ansatz}\\
	\vspace*{1cm}
	%\includegraphics[height=6cm]{#1}
	\vfill
	\normalsize
	%\newcolumntype{x}[1]{>{\raggedleft\arraybackslash\hspace{0pt}}p{#1}}
	\begin{tabular}{rl}%{6cm}p{7.5cm}}
		\rule{0mm}{5ex}\textbf{Autoren:} & \hspace*{-0.5em}\begin{tabular}[t]{r}#3\end{tabular} \\ 
		\rule{0mm}{5ex}\textbf{Leiter:} & Prof. Dr. rer. nat. Carsten Kern \\ 
		\rule{0mm}{5ex}\textbf{Abgabedatum:} & #4 \\ 
	\end{tabular} 
\end{center}
\pagebreak
}

% ----------------------------------------------------------------------------------------------------------
% Titelseite
% ----------------------------------------------------------------------------------------------------------
\MyTitlepage{}{
\texttt{simon1.hofmeister@st.oth-regensburg.de}\\
\texttt{naddia1.matsko@st.oth-regensburg.de}\\
\texttt{monika.silber@st.oth-regensburg.de}\\
\texttt{simon.wasserburger@st.oth-regensburg.de}}
{15.03.\the\year}

\setcounter{page}{1} 
% ----------------------------------------------------------------------------------------------------------
% Inhaltsverzeichnis
% ----------------------------------------------------------------------------------------------------------
\tableofcontents
\pagebreak


% ----------------------------------------------------------------------------------------------------------
% Inhalt
% ----------------------------------------------------------------------------------------------------------
% Abstände Überschrift
\titlespacing{\section}{0pt}{12pt plus 4pt minus 2pt}{-6pt plus 2pt minus 2pt}
\titlespacing{\subsection}{0pt}{12pt plus 4pt minus 2pt}{-6pt plus 2pt minus 2pt}
\titlespacing{\subsubsection}{0pt}{12pt plus 4pt minus 2pt}{-6pt plus 2pt minus 2pt}

% Kopfzeile
\renewcommand{\sectionmark}[1]{\markright{#1}}
\renewcommand{\subsectionmark}[1]{}
\renewcommand{\subsubsectionmark}[1]{}
\lhead{Kapitel \thesection}
\rhead{\rightmark}

\onehalfspacing
\renewcommand{\thesection}{\arabic{section}}
\renewcommand{\theHsection}{\arabic{section}}
\setcounter{section}{0}
\pagenumbering{arabic}
\setcounter{page}{1}

% ----------------------------------------------------------------------------------
% Kapitel: Einleitung
% ----------------------------------------------------------------------------------
\section{Einleitung}

\newpage

% ----------------------------------------------------------------------------------
% Kapitel: Related Work/Theorie
% ----------------------------------------------------------------------------------
\section{Related Work}

\subsection{Monte Carlo Tree Search}

\newpage

% ----------------------------------------------------------------------------------
% Kapitel: Implementierung/Umsetzung
% ----------------------------------------------------------------------------------
\section{Implementierung}

\subsection{Monte Carlo Tree Search}
Um den Monte Carlo Tree Search (MCTS) für Reversi zu realisieren wurden die Klassen MCTS und Node angelegt. Letztere enhält zwei überladene Konstruktoren zum Anlegen von Wurzel- und Kindknoten. Ein Node enthält zusätzliche Attribute. Sowohl der Elternknoten als auch eine ArrayList vom Typ Node, die die direkten Kinder enthält, werden abgespeichert. Die Anzahl, wie oft ein Knoten besucht wurde, wird in der Variable numVisited hinterlegt. Außerdem gespeichert wird das Ergebnis eines simulierten Spieldurchlaufs in simulationReward. Da der aktuelle Spielzustand durch das Environment, das ebenfalls die derzeitige Repräsentation des Playgrounds beinhaltet, definiert ist, wird dies in Node hinterlegt. Des Weiteren wird in dem Attribut nextPlayer festgehalten, welcher Spieler als Nächstes an der Reihe ist.
Bei der Instanziierung durch die Konstruktoren werden sinnvolle initiale Werte vergeben. numVisited und simulationReward werden auf 0 beziehungsweise 0.0 gesetzt. Für den Wurzelknoten gilt, dass er keinen Parent besitzt, für alle weiteren Knoten wird der Parent übergeben und gesetzt. Die Children werden zunächst durch eine leere Liste initialisiert. Der nextPlayer wird ebenfalls übergeben und gesetzt. Außerdem zu erwähnen ist die Methode calculateUCT(), die den Upper Confidence Bound applied to Trees (UCT) für einen Knoten berechnet. Sie ermittelt zunächst die Exploitation-Komponente, indem sie den simulationReward durch die Anzahl an Besuchen dividiert. Zur Berechnung der Exploration wird die Wurzel aus dem Wert, wie oft der Parent besucht wurde geteilt durch den Wert die oft der aktuelle Knoten besucht wurde +1, gezogen. Um den Kompromiss zwischen diesen beiden Komponenten zu kontrollieren, wird die Exploration mit der A-priori-Wahrscheinlichkeit für einen Zug multipliziert. Diese wird im neuralen Netz trainiert.
Hinsichtlich der Klasse MCTS ist festzuhalten, dass diese im der Klasse Agent über den Konstruktoraufruf instanziiert wird. Dieser verlangt das Environment und den Player als Übergabeparameter und legt daraufhin einen neuen Wurzelknoten, sowie eine leere ArrayList vom Typ Node, die die Blattknoten beinhaltet, die im späteren Verlauf simuliert werden müssen. Für die Simulation muss beachtet werden, dass das Environment geklont und somit eine tiefe Kopie erzeugt werden muss, damit der tatsächliche Spielzustand nicht unbeabsichtigt manipuliert wird. Dabei ist festzuhalten, dass dies zweimal stattfindet. Einmal, wenn ein neuer Baum aufgebaut wird, somit erhält der neue Wurzelknoten und ebenfalls jedes seiner Kinder jeweils einen eigenen Klon. Für die Kinderknoten gilt, dass diese ihre Environment-Instanz an ihre Kinder weitergeben und diese somit innerhalb derselben Instanz agieren. Um den MCTS zu starten, wird die Methode searchBestTurn() aufgerufen. Diese expandiert zunächst den Wurzelknoten, indem sie alle im aktuellen Zustand möglichen validen Züge ermittelt und durch diese iteriert. Die Methode getPossibleTurns() gibt diese zurück. Sie iteriert über das gesamte Spielfeld und prüft dabei mithilfe der Methode validateTurnPhase1 im Environment, an welcher Stelle ein gültiger Zug gemacht werden kann. Für jeden dieser Züge wird der Kinderknoten angelegt, sowie als unbesuchte Blattknoten abspeichert. Außerdem wird ermittelt, welcher Spieler als Nächster einen Zug machen darf und der simulierte derzeitige Spielzustand anhand des Zuges aktualisiert. Daraufhin werden die unbesuchten Blattknoten, die am Anfang den Kinderknoten der Wurzel entsprechen, durchlaufen. Dabei wird in jedem dieser eine Simulation gestartet, die einen Spielverlauf bis zum Spielende anhand zufällig ausgewählter möglicher Züge durchspielt. Mithilfe der rekursiven Funktion simulate() erfolgt die Simulation eines Spiels, sowie die Backpropagation des Ergebnisses bei Spielende. Sie enthält letztendlich den Reward als Rückgabewert. Ein Zufallszug wird durch eine Instanz der Java-Klasse Random generiert. Hierbei wird ein zufälliger Integer erzeugt, der durch eine Modulo-Operation auf den Größenbereich abgestimmt wird, der dem der Anzahl der möglichen Züge entspricht. Die resultierende Zahl gibt den auszuwählenden Zug innerhalb der ArrayList an. Für einen neu gewählten Zug wird ein neuer Knoten angelegt und dessen Variablen, die zur Berechnung des UCT relevant sind, aktualisiert. Der Reward wird durch den rekursiven Aufruf der Funktion erhöht. Die Anzahl an Besuchen wird jeweils um den Wert eins erhöht. Wenn keine weiteren validen Folgezüge ermittelt werden können, bedeutet das das Spielende und der Reward für die Spielausgang wird anhand der Funktion rewardGameState() berechnet. Diese erhält als Parameter das Environment, sowie den Spieler, für den der Reward kalkuliert werden soll. Indem der gesamte Playground durchlaufen und gezählt wird, wie viele Steine vom übergebenen Spieler enthalten sind, errechnet sich die Bewertung des Spiels. Abschließend werden die Werte für Anzahl Besuche und Reward ebenfalls im Wurzelknoten aktualisiert.


\newpage

% ----------------------------------------------------------------------------------
% Kapitel: Allgemeine Informationen/Organisation
% ----------------------------------------------------------------------------------
\section{Organisation}

\subsection{Team und Aufgabenverteilung}
% Beschreiben Sie in diesem Abschnitt Ihr Team. Welche Person hat welche Aufgaben wahrgenommen, wie wurden
% Aufgaben aufgeteilt und wie wurde kommuniziert, etc.

\subsection{Kommunikation}

\subsection{Versionskontrolle}

\subsection{OS, IDE und Programmiersprache}

\subsection{Testumgebungen}

\subsection{Projekt-Dokumentation}

\newpage
% ----------------------------------------------------------------------------------
% Kapitel: Fazit
% ----------------------------------------------------------------------------------
\section{Fazit}
% Beschreiben Sie in diesem Abschnitt u.a.\ was Ihnen an diesem Fach gefallen hat und welche
% Verbesserungsvorschläge Sie für künftige Veranstaltungen haben. Was konnten Sie dazulernen, in welchen
% Bereichen haben Sie sich verbessert. Welche Problemsituationen gab es während der Projekterstellung, wie
% sind Sie diese angegangen und wie haben Sie diese gelöst. Was haben Sie evtl.\ vermisst.


%\newpage
% ----------------------------------------------------------------------------------
% Kleine Einführung in LaTeX-Elemente
% ----------------------------------------------------------------------------------
% \section{\LaTeX-Elemente}
% Dieser Abschnitt soll nicht Bestandteil des Projektberichtes sein, sondern beinhaltet lediglich einige
% Informationen über \LaTeX-Distributionen, Editoren und \LaTeX-Elemente, die Ihnen beim Einstieg in das
% \LaTeX-Textsatzsystem helfen sollen.

% \subsection{\LaTeX-Distributionen nach Betriebssystemen}

% \subsubsection{\LaTeX-Distributionen}

% Folgende Haupt-\LaTeX-Distributionen stehen Ihnen zur Verfügung:
% \begin{itemize}
%  \item Windows:\quad \texttt{MiKTeX}\quad Webseite:\quad\url{http://www.miktex.org}
%  \item Linux/Unix:\quad \texttt{TeX Live}\quad Webseite:\quad\url{http://tug.org/texlive/}
%  \item Mac OS:\quad \texttt{MacTeX}\quad Webseite:\quad\url{http://www.tug.org/mactex/}
% \end{itemize}

% \subsubsection{\LaTeX-Editoren}
% Auf folgenden Webseiten können Sie einige hilfreiche \LaTeX-Editoren finden:
% \begin{itemize}
%  \item Windows/Linux/Mac OS: \url{http://www.xm1math.net/texmaker/}
%  \item Windiws: \url{http://www.texniccenter.org/}
%  \item Mac OS: \url{http://pages.uoregon.edu/koch/texshop/}
% \end{itemize}

%Falls bei den oben genannten Editoren kein passender vorhanden war, findet sich auf Wikipedia eine  %Zusammenstellung vieler weiterer \LaTeX-Editoren:\\[1em]
%\hspace*{3cm}\url{https://en.wikipedia.org/wiki/Comparison_of_TeX_editors}


%\subsection{Unterabschnitt}
%Zum Einfügen eines Bildes, siehe Abbildung \ref{fig:reversi01}, wird die
% \textit{minipage}-Umgebung 
%genutzt, da die Bilder so gut positioniert werden können.

%\vspace{1em}
%\begin{minipage}{\linewidth}
%	\centering
%	\includegraphics[width=0.6\linewidth]{pics/gamefield01.png}
%	\captionof{figure}[Spielfeld 01]{Unbespieltes Spielfeld\footnotemark }
%	\label{fig:reversi01}
%\end{minipage}
%\footnotetext{Diesem Spielfeld wurden noch keine Spieler zugewiesen (daher die
% dunklen Spielsteine)}

%Nachdem das Spiel gestartet wurde und beiden Spielphasen durchlaufen wurden, siegt
% schließlich der %Spieler mit der Farbe rot.

%\vspace{1em}
%\begin{minipage}{\linewidth}
%	\centering
%	\includegraphics[width=0.6\linewidth]{pics/gamefield02.png}
%	\captionof{figure}[Spielfeld 02]{Finales Spielfeld\footnotemark }
%	\label{fig:reversi2}
%\end{minipage}
%\footnotetext{Das Spielfeld nach der Zug- und Bombenphase. Spieler rot gewinnt
% eindeutig.}

%\subsection{Tabellen}
%In diesem Abschnitt wird eine Tabelle (siehe Tabelle \ref{tab:beispiel}) dargestellt.

%\vspace{1em}
%\begin{table}[!h]
%	\centering
%	\begin{tabular}{|l|l|l|}
%		\hline
%		\textbf{Name} & \textbf{Name} & \textbf{Name}\\
%		\hline
%		1 & 2 & 3\\
%		\hline
%		4 & 5 & 6\\
%		\hline
%		7 & 8 & 9\\
%		\hline
%	\end{tabular}
%	\caption{Beispieltabelle}
%	\label{tab:beispiel}
%\end{table}


%\subsection{Auflistung}
%Für Auflistungen wird die \textit{enumerate}- oder \textit{itemize}-Umgebung genutzt.
%
%\begin{itemize}
%	\item Nur
%	\item ein
%	\item Beispiel.
%\end{itemize}
%
%\subsection{Listings}
%Zuletzt ein Beispiel für ein Listing, in dem Quellcode eingebunden werden kann, siehe Listing \ref{lst:arduino}.
%
%\vspace{1em}
%\begin{lstlisting}[caption=Arduino Beispielprogramm, label=lst:arduino]
%int ledPin = 13;
%void setup() {
%    pinMode(ledPin, OUTPUT);
%}
%void loop() {
%    digitalWrite(ledPin, HIGH);
%    delay(500);
%    digitalWrite(ledPin, LOW);
%    delay(500);
%}
%\end{lstlisting}
%
%\subsection{Tipps}
%Die Quellen befinden sich in der Datei \textit{quellen.bib}. Eine Buch- und eine Online-Quelle sind beispielhaft eingefügt. [Vgl. \cite{buch}, \cite{online}]
%
%\pagebreak

% ----------------------------------------------------------------------------------------------------------
% Literatur
% ----------------------------------------------------------------------------------------------------------
\renewcommand\refname{Quellenverzeichnis}
\bibliographystyle{alpha}
\bibliography{quellen}
\pagebreak

% ----------------------------------------------------------------------------------------------------------
% Anhang
% ----------------------------------------------------------------------------------------------------------
\pagenumbering{Roman}
\setcounter{page}{1}
%\lhead{Anhang \thesection}

\begin{appendix}
\section*{Anhang}
%\phantomsection
\addcontentsline{toc}{section}{Anhang}
\addtocontents{toc}{\vspace{-0.5em}}

\end{appendix}

\end{document}